\begin{textsample}{POS Dim 3 – persona\_gpt – Score 34.00 – t954\_gpt.txt}  \label{ex:f3_pos_022}
Our food \textbf{system} is failing us.
From the \textbf{soil} beneath our feet to the animals in \textbf{factory} farms, from the \textbf{health} of farmers to the food on our \textbf{plates}, the industrial \textbf{model} is built on \textbf{chemicals}, \textbf{waste}, and short-term profit rather than care for people and the planet.
The \textbf{scale} of the damage is staggering.
In 2007, 269 \textbf{tonnes} of \textbf{pesticides} were \textbf{used} globally every hour.
That \textbf{means} toxic \textbf{chemicals} being sprayed into the air we breathe, the water we drink, and the ecosystems that support life.
This is not a marginal issue at the edges of \textbf{agriculture} ; it is the operating \textbf{system} of industrial food \textbf{production}.
At the same \textbf{time}, this \textbf{system} is astonishingly wasteful.
One-third of all food produced globally is \textbf{wasted}.
That is not just a statistic about rotting leftovers ; it represents land, water, \textbf{energy}, and labour thrown away.
The \textbf{amount} of food \textbf{wasted} each year could feed 149 million people in Europe.
While millions go hungry, mountains of edible food are discarded along the \textbf{supply} \textbf{chain} and in our homes.
The hidden \textbf{cost} of our \textbf{diets} is also \textbf{measured} in water.
For just one person ’s annual \textbf{consumption} of \textbf{meat} and \textbf{dairy}, an estimated 403,000 litres of water are \textbf{used}.
Behind every steak, every glass of milk, every slice of cheese lies a vast, often invisible drain on rivers, aquifers, and freshwater ecosystems.
This is the reality of a \textbf{broken} industrial food \textbf{system} : chemical-intensive, resource-hungry, and indifferent to the wellbeing of farmers, communities, animals, and nature.
But this \textbf{system} is not inevitable, and it is not unchangeable.
Ecological farming offers a different path.
It puts people and \textbf{health} at the centre, rather than \textbf{chemicals} and corporate \textbf{control}.
It values biodiversity instead of monocultures, and it respects animals as living beings rather than treating them as \textbf{units} in a \textbf{production} line.
Ecological farming integrates rigorous science with local knowledge, recognising that farmers and communities are not just end-users of \textbf{technology}, but co-creators of \textbf{solutions}.
This is not simply about changing how food is grown ; it is about changing who the food \textbf{system} serves.
Individual choices can help drive that \textbf{shift}.
Every \textbf{time} we buy ecological, local, and seasonal food, we support farming that works with nature instead of against it.
When we choose to shop at farmers ’ \textbf{markets}, we strengthen direct relationships between growers and eaters, keeping more value in local communities and reducing the distance food travels.
When we eat less \textbf{meat}, we reduce pressure on land, water, and climate, and open space for more diverse, sustainable farming \textbf{systems}.
These actions \textbf{may} seem small, but together they build a people-powered movement.
When enough of us demand ecological food and farming, \textbf{markets} respond.
When \textbf{markets} \textbf{shift}, governments follow with policies and support that can \textbf{scale} up ecological practices and phase out chemical-intensive industrial \textbf{agriculture}.
The \textbf{broken} industrial food \textbf{system} is not just a policy \textbf{problem} or a scientific challenge ; it is a democratic issue.
Who decides what we grow, how we grow it, and who benefits?
By changing what we put on our \textbf{plates} and where we spend our \textbf{money}, we reclaim some of that power.
The \textbf{numbers} on \textbf{pesticides}, food \textbf{waste}, and water \textbf{use} make it clear that the status quo is unsustainable.
The promise of ecological farming shows that another future is possible—one where our food \textbf{system} nourishes people, protects farmers, restores biodiversity, and respects animals.
That future \textbf{will} not arrive on its own.
It \textbf{will} be built, bite by bite and choice by choice, by all of us.

% matched lemmas: agriculture, amount, break, chain, chemical, consumption, control, cost, dairy, diet, energy, factory, health, market, may, mean, measure, meat, model, money, number, pesticide, plate, problem, production, scale, shift, soil, solution, supply, system, technology, time, tonne, unit, use, waste, will
\end{textsample}
